% Pursuing Stacks (cotes 134-2 to 134-8): archive cote 134-4, site page 28 (PDF page 29), typescript page 116.
% typescript. Draft compared: ps3.tex:1377-1443.
% AI-transcribed by Claude Opus 5.5 (claude-opus-5-5, adaptive thinking, effort high), Message Batches API, batch ps2 (msgbatch_01NVwr1Yuo8qJf5rykeGBXPK), 2026-11-12; not yet checked by eye.
\scanpage{134-4/28}
\mspage{116}
that the product of a finite family of contractible objects of M is again
contractible.
\inserted{We assume M stable under finite products.}
\margin{X} C) \inserted[typed]{\uline{Weak}} \uline{Homotopy intervals}. The type  of structure we've in view is a
\xout{\unsure{sub}}set  of intervals in M
\[ \hand{\Sigma}_h \subset \mathrm{Int}(M), \]
called the "weak homotopy intervals". No basic assumption on this set,
it seems;  the natural ``\uline{saturation} \uline{condition}'' is the following:
(Sat 3)  \xout{\unsure{Any}} Any interval \inserted[typed]{$\underline{I}$ =} $(I,\hand{\delta}_0,\hand{\delta}_1)$ in M,  such that the sections $\hand{\delta}_0,\hand{\delta}_1$
of I be $\hand{\Sigma}_h$-homotopic (see below),  is in $\hand{\Sigma}_h$ .

THe assumption on $\underline{I}$ means,  explicitly,  that there exists a finite
chain of sections of I
\[ s_0=\hand{\delta}_0, s_1, \ldots , s_N=\hand{\delta}_1 , \]
joining $\hand{\delta}_0$ to $\hand{\delta}_1$ ,  and for two consecutive  $s_i, s_{i+1}$  an interval
$\underline{J}$ in $\hand{\Sigma}_h$ ,  and a \xout{map} map of intervals \inserted[typed]{from J or} $\check{\underline{J}}$ to $(I,s_i,s_{i+1})$  i.e. a map
$J \to I$ ,  mapping the two \xout{\illegible{1}} \inserted[typed]{given} sections of J  one into $s_i$ , the other
into $s_{i+1}$ (without specification which is mapped into which).

THe significance of this saturation condition becomes clear in terms
of the second dotted arrow of the organigram.  Namely,  in terms of any
subset $\hand{\Sigma}_h$ of \xout{the} Int(M),  we get a corresponding homotopy relation bet-%
ween maps,  say $R_h$ ,  which is the \xout{\illegible{1}} equivalence relation \xout{\illegible{1}} in
Fl(M)  generated by the"elementary" homotopy relation \xout{\illegible{1}} (with
redpect to $\hand{\Sigma}_h$) between maps f,g in M,  namely the\xout{t} relation  $R_0$
\[ f \underset{R_0}{\sim} g \overset{\mathrm{def}}{\Longleftrightarrow} \text{there exists $\underline{I}$ in $\hand{\Sigma}_h$ ,  and an $\underline{I}$-homotopy} \]
\[ \text{from f to g .} \]
THe corresponding equivalence relation $R_h$ in Fl(M) is already saturated
namely stable under compositions,  moreover it satisfies condition $C_{12}$
above  -  thus we may view this homotopy relation as defined in terms of
a homotopism\xout{s} notion  -  thus in fact the second dotted arrow should go
from  3)  to 2)  rather than from 3) to 1) !  Now,  if \xout{it is clear that}we
we \xout{l}look at the subset $\overline{\hand{\Sigma}}_h$  of Int(M)  defined in terms of $\over{\unsure{\Sigma}}{R}_h$  as in A)
(namely the set of ``\inserted{weak} \unsure{h|H}omotopy intervals  with respect to $R_h$''),  we get \xout{x}
\[ R_h \subset \overline{R}_h \quad , \]
and the equality holds iff $R_h$ satisfies (Sat 3)?  At the same time,  in
case of arbitrary $R_h$ ,  we get the construction of it's saturation, \inserted{$\overline{R_h}$} which
\xout{\illegible{1}}may of course  be described alternatively as the smallest \xout{subset} saturated
subset of Int(M) containing $R_h$ .

We'll call \uline{weak homotopy interval structure} on M,  \xout{the} any \deleted{\unsure{sub}}set
\margin{28}
